\chapter{Post-Traumatic Stress}

\section*{Definition}
Post-traumatic stress disorder (PTSD) is a syndrome that develops after exposure to actual or threatened death, serious injury, or sexual violence, characterized by intrusion symptoms, avoidance, negative alterations in cognition and mood, and marked changes in arousal and reactivity present for more than 1 month.

\section*{What Happens in Your Body}
PTSD involves hyperactivation of the amygdala with impaired top-down regulation from the ventromedial prefrontal cortex and hippocampus. The HPA axis shows altered cortisol regulation (often low cortisol with increased CRH), while the locus coeruleus drives noradrenergic hyperarousal. Dysfunctional fear extinction learning prevents the resolution of conditioned fear responses.

\section*{Causes}
\begin{itemize}
  \item \textbf{Traumatic events:} Combat, sexual assault, physical assault, childhood abuse, motor vehicle accidents, natural disasters, terrorist attacks, life-threatening medical illness
  \item \textbf{Psychological:} Perceived threat to life, peritraumatic dissociation, low social support, chronic stress, prior trauma history
  \item \textbf{Genetic:} Heritability 30--40\%; polymorphisms in FKBP5, CRHR1, and SLC6A4 (serotonin transporter) genes
  \item \textbf{Neurobiological:} Smaller hippocampal volume, reduced prefrontal cortical thickness
  \item \textbf{Risk factors:} Female sex, younger age, lower socioeconomic status, preexisting anxiety/depression, prior PTSD
\end{itemize}

\section*{When to See a Doctor}
See your doctor if:
\begin{itemize}
  \item Intrusive memories, nightmares, or flashbacks persist beyond 1 month after trauma
  \item Avoidance behaviors significantly limit functioning
  \item There is suicidal ideation, severe hyperarousal, or substance use as a coping mechanism
\end{itemize}

\section*{Related Symptoms}
\begin{itemize}
  \item Recurrent intrusive memories, nightmares, flashbacks, hypervigilance, exaggerated startle response, emotional numbness, avoidance of trauma reminders, negative beliefs, irritability, dissociation, depression, substance use
\end{itemize}
\textbf{Important:} PTSD affects approximately 7--8\% of the population at some point in their lives, with women twice as likely as men to receive the diagnosis.

\section*{Self-Care / Home Management}
\begin{itemize}
  \item Establish a consistent daily routine and sleep schedule
  \item Limit exposure to trauma reminders (news, social media) and use grounding when exposed
  \item Engage in regular physical activity to reduce hyperarousal
  \item Avoid alcohol and other substances that interfere with trauma processing
\end{itemize}

\section*{How Your Doctor Will Evaluate You}
\begin{itemize}
  \item Clinical interview using DSM-5 (PTSD) or DSM-5-TR criteria
  \item Structured instruments: Clinician-Administered PTSD Scale for DSM-5 (CAPS-5); PCL-5 for screening
  \item Assess for complex PTSD (ICD-11) if early childhood trauma, dissociation, affect dysregulation
  \item Screen for comorbid depression, substance use, suicidal ideation, and sleep disturbance
\end{itemize}

\section*{Treatment Options}
\begin{itemize}
  \item First-line trauma-focused psychotherapy: cognitive processing therapy (CPT), prolonged exposure (PE), or EMDR
  \item First-line pharmacotherapy: SSRIs (sertraline, paroxetine) or SNRI (venlafaxine); fluoxetine is also effective
  \item Adjunctive prazosin for trauma-related nightmares
  \item Stabilization skills (grounding, distress tolerance) before trauma processing for complex PTSD
  \item Combination therapy (CBT plus medication) for severe or treatment-resistant cases
\end{itemize}

\clearpage
